% =====================================================================
% 08_validation.tex — démonstration, questions-réponses et validation
% =====================================================================
\chapter{Démonstration, questions-réponses et validation}
\label{chap:validation}

\section{Objectif de la validation}

Ce chapitre prépare la soutenance. Il ne présente pas de nouveau mécanisme :
il organise les preuves, propose un déroulé de dix minutes et rassemble les
questions les plus probables. L'objectif n'est pas de tout montrer, mais de
choisir quelques démonstrations parlantes, de savoir expliquer les choix et de
réagir proprement à une erreur volontaire.

\begin{objectifbox}[Principe directeur]
Une soutenance réussie montre trois choses : le projet fonctionne, les choix
sont compris, et les erreurs affichées sont maîtrisées. Il vaut mieux montrer
cinq preuves solides et expliquer les limites que tenter de tout exécuter et
perdre le fil.
\end{objectifbox}

\section{Script de démonstration de dix minutes}

Le script \code{demo\_soutenance.sql} reprend les blocs dans l'ordre. Tous les
mouvements sont annulés par \code{ROLLBACK} ; l'état final reste 6 clients,
6 comptes courants, 4 comptes d'épargne, 3 comptes à terme, 15 opérations,
9 signataires, 2 alertes et 9 affectations. Le script a été testé : zéro objet
invalide.

\begin{table}[htbp]
  \centering
  \small
  \caption{Déroulé de la démonstration de dix minutes.}
  \label{tab:demo-10min}
  \begin{tabularx}{\textwidth}{@{}C{1.8cm}L{4.8cm}X@{}}
    \toprule
    \entete{Temps} & \entete{Séquence} & \entete{Preuve à montrer} \\
    \midrule
    0:00--0:40 & Introduction & Utilisateur \code{MASTER1}, zéro objet
    invalide, script rejouable. \\
    0:40--2:00 & Partie A : schéma & Types, tables, comptes, hiérarchie,
    \code{VARRAY} et \emph{nested table}. \\
    2:00--3:00 & Partie B : manipulations & Téléphone du client 3, contrainte
    différée, \code{drop\_client} en sortie. \\
    3:00--4:15 & Partie C : requêtes & 6.10 (CC2 négatif), 6.20 (signataires
    de CC7). \\
    4:15--6:00 & Partie D : méthodes et transactions & \code{nbCepargne},
    \code{soldeTotal}, intérêts, \code{passeMvt} avec rollback. \\
    6:00--7:30 & Partie D : triggers & Deux refus, un succès avec rollback,
    explication du \code{COMPOUND}. \\
    7:30--8:45 & Partie E : vue objet & \code{Agence\_VOR}, collections,
    parrainage, mutation en sortie. \\
    8:45--10:00 & Conclusion & Comptage des objets, zéro invalide, forces et
    limites. \\
    \bottomrule
  \end{tabularx}
\end{table}

\subsection{Introduction et contrôle de l'état}

La première commande vérifie l'utilisateur ; la seconde prouve qu'aucun objet
n'est invalide. C'est la preuve la plus importante, car elle condense la
réussite du déploiement.

\begin{lstlisting}[caption={Contrôle d'entrée en soutenance.},label={lst:demo-intro}]
SELECT USER FROM dual;

SELECT COUNT(*) AS objets_invalides
  FROM user_objects
 WHERE status <> 'VALID';
\end{lstlisting}

Le résultat attendu est \code{MASTER1}, puis \code{0}. Il faut le dire avant
de commenter les types : cela coupe court aux doutes sur la compilation.

\subsection{Partie A : schéma}

Les requêtes de contrôle listent les types, les tables et les comptes. Il ne
faut pas afficher toutes les lignes : se contenter des premières et expliquer
la structure. Les 19 types métier, les 11 tables logiques et la table de
stockage des signataires suffisent à situer le projet.

\begin{lstlisting}[caption={Contrôle du schéma en démonstration.},label={lst:demo-schema}]
SELECT type_name, typecode FROM user_types ORDER BY type_name;
SELECT table_name FROM user_tables ORDER BY table_name;
SELECT COUNT(*) AS clients FROM client;
SELECT COUNT(*) AS comptes_courants FROM cpt_courant;
\end{lstlisting}

Le commentaire oral doit insister sur trois points : \code{compte\_type} est
non instanciable, il existe une table par sous-type, et \code{client} conserve
un OID système pour que les \code{REF} puissent entrer dans la \emph{nested
table} des signataires.

\subsection{Partie B : manipulations}

La vérification du téléphone du client 3 montre l'effet de
\code{change\_Portable} : le deuxième téléphone doit être
\code{0777-11-22-33}. Pour \code{drop\_client}, il vaut mieux montrer la sortie
de \code{\_FULL.out} que supprimer un client réel pendant la démonstration.

\begin{lstlisting}[caption={Vérification du téléphone du client 3.},label={lst:demo-b5e}]
SELECT c.num_client, t.COLUMN_VALUE AS telephone
  FROM client c, TABLE(c.telephones) t
 WHERE c.num_client = 3;
\end{lstlisting}

\subsection{Partie C : deux requêtes suffisent}

La requête 6.10 montre le découvert de CC2 et le déréférencement du conseiller.
La requête 6.20 montre la jointure entre une \emph{nested table} et la table
\code{client} via \code{s.ref\_client = REF(c)}. Ces deux requêtes couvrent le
déréférencement, les collections et les jointures.

\begin{lstlisting}[caption={Requêtes 6.10 et 6.20 en démonstration.},label={lst:demo-c}]
SELECT cc.num_compte, cc.solde, c.nom, k.nom AS conseiller
  FROM cpt_courant cc, client c, conseiller k
 WHERE cc.ref_client = REF(c)
   AND c.ref_conseiller = REF(k)
   AND cc.solde < 0;

SELECT c.num_client, c.nom,
       c.adresse.num_rue||' '||c.adresse.rue||', '||c.adresse.ville AS adresse,
       s.droit
  FROM cpt_courant cc, TABLE(cc.signataires) s, client c
 WHERE cc.num_compte = 'CC7'
   AND s.ref_client = REF(c)
 ORDER BY c.num_client;
\end{lstlisting}

\begin{attentionbox}[Ne pas montrer 6.24 après la partie D]
L'opération de 1\,500\,000~DA a été supprimée avec CC1. Si la question porte
sur 6.24, montrer \code{\_C.out} ou la section correspondante de
\code{\_FULL.out}. Il ne faut pas laisser croire que la requête ne renvoie
rien.
\end{attentionbox}

\subsection{Partie D : méthodes, transactions et triggers}

Les méthodes de \code{client} et de \code{cpt\_epargne} se montrent en une
requête. Ensuite, \code{passeMvt} est exécuté dans un bloc qui affiche le solde
avant et après, puis annule par \code{ROLLBACK}. Le solde de CC7 revient à
185\,000, ce qui prouve l'absence de \code{COMMIT} interne.

\begin{lstlisting}[caption={Méthodes de client et d'épargne.},label={lst:demo-methodes}]
SELECT c.num_client, c.nom, c.nbCepargne() AS nb_epargne,
       c.soldeTotal() AS solde_total
  FROM client c
 ORDER BY c.num_client;

SELECT ce.num_compte, ce.solde, ce.calculeInterets() AS a_l_echeance,
       ce.sortieAnticipee() AS sortie_anticipee
  FROM cpt_epargne ce
 ORDER BY ce.num_compte;
\end{lstlisting}

\begin{lstlisting}[caption={\code{passeMvt} avec rollback.},label={lst:demo-passeMvt}]
DECLARE
  v_cc    cptCourant_type;
  v_avant NUMBER;
  v_apres NUMBER;
BEGIN
  SELECT solde INTO v_avant FROM cpt_courant WHERE num_compte = 'CC7';
  DBMS_OUTPUT.PUT_LINE('Avant : '||v_avant);
  SELECT VALUE(cc) INTO v_cc FROM cpt_courant cc WHERE cc.num_compte = 'CC7';
  v_cc.passeMvt(4, 10000);
  SELECT solde INTO v_apres FROM cpt_courant WHERE num_compte = 'CC7';
  DBMS_OUTPUT.PUT_LINE('Apres : '||v_apres);
  ROLLBACK;
END;
/
\end{lstlisting}

Deux refus de trigger suffisent : un compte bloqué et un dépassement de
plafond. Ils doivent être capturés par \code{EXCEPTION WHEN OTHERS} et affichés
comme « refus attendu ». Un succès avec \code{ROLLBACK} montre ensuite que le
trigger \code{AFTER} a bien mis à jour le solde, sans persister la modification.

\begin{lstlisting}[caption={Refus d'une opération sur un compte bloqué.},label={lst:demo-refus}]
BEGIN
  INSERT INTO operation (id_op, ref_client, ref_compte, date_op,
                         montant, type_op, commission)
  VALUES (seq_operation.NEXTVAL,
          (SELECT REF(c) FROM client c WHERE c.num_client = 5),
          (SELECT REF(cc) FROM cpt_courant cc WHERE cc.num_compte = 'CC5'),
          TRUNC(SYSDATE), -1000, 'RETRAIT', 0);
  DBMS_OUTPUT.PUT_LINE('Inattendu : accepte');
EXCEPTION
  WHEN OTHERS THEN
    DBMS_OUTPUT.PUT_LINE('Refus attendu : '||SQLERRM);
END;
/
\end{lstlisting}

\subsection{Partie E : vue objet}

Les trois requêtes sur \code{Agence\_VOR} montrent la vue, les comptes et les
conseillers. La mutation est plus délicate : elle modifie l'historique. En
démonstration, il est préférable de montrer la sortie déjà exécutée de
\code{\_FULL.out} plutôt que de modifier l'état en direct. Si le jury insiste,
expliquer que la méthode ferme l'affectation courante et en ouvre une nouvelle,
puis signaler que le responsable doit être revu.

\begin{lstlisting}[caption={Interrogation de la vue objet.},label={lst:demo-e}]
SELECT code, nom, ville FROM Agence_VOR ORDER BY code;

SELECT v.code, c.num_compte, c.solde
  FROM Agence_VOR v, TABLE(v.comptes) c
 ORDER BY v.code, c.num_compte;

SELECT v.code, k.nom, k.prenom
  FROM Agence_VOR v, TABLE(v.conseillers) k
 ORDER BY v.code, k.nom;
\end{lstlisting}

\subsection{Conclusion de la démonstration}

Le comptage final des objets et le contrôle des objets invalides ferment la
démonstration. Il faut commenter les forces — héritage, collections, triggers
de sécurité, vue objet — et les limites — opérations sur les comptes courants,
virement CC vers CC, pas de remplacement automatique du responsable, pas
d'interface graphique, pas de tests de concurrence.

\begin{lstlisting}[caption={Contrôle final des objets.},label={lst:demo-fin}]
SELECT object_type, COUNT(*) AS nb
  FROM user_objects
 GROUP BY object_type
 ORDER BY object_type;

SELECT COUNT(*) AS objets_invalides
  FROM user_objects
 WHERE status <> 'VALID';
\end{lstlisting}

\section{Questions-réponses}

Les quarante-cinq questions ci-dessous couvrent la conception, les parties A à
E et la validation. Chaque réponse est donnée en deux temps : une réponse
courte, puis un complément si le jury insiste.

\subsection{Conception et modèle}

\qa{1}{Pourquoi des tables objet et pas uniquement du relationnel ?}
{Le sujet impose types objet, héritage, \code{REF}, collections, méthodes,
type récursif et vue objet.}
{Le relationnel pur ne permet pas de démontrer ces mécanismes. On garde du
relationnel pour l'historique, les alertes et l'audit, là où c'est plus simple.}

\qa{2}{Pourquoi \code{compte\_type} est-il \code{NOT FINAL}
\code{NOT INSTANTIABLE} ?}
{\code{NOT FINAL} autorise l'héritage ; \code{NOT INSTANTIABLE} interdit un
compte générique.}
{La règle métier dit que la banque ne gère jamais de compte « nu ». Chaque
compte est obligatoirement courant, épargne ou terme.}

\qa{3}{Pourquoi une table par sous-type et pas une seule table
\code{compte} ?}
{Chaque sous-type a ses attributs propres ; cela évite les colonnes nullables
et respecte l'interdiction du compte générique.}
{Une table unique aurait un discriminant et des colonnes inutiles selon le
type. La hiérarchie objet est plus propre et permet le polymorphisme de
\code{fraisTenue}.}

\qa{4}{Pourquoi \code{client\_type} n'a-t-il pas d'attribut
\code{agence} ?}
{Parce qu'un client est suivi par un conseiller, et qu'un conseiller peut
changer d'agence. L'agence se déduit de l'affectation du conseiller.}
{Mettre \code{ref\_agence} dans le client dupliquerait l'information et
poserait des problèmes de mise à jour.}

\qa{5}{Pourquoi \code{affectation} est-elle une table relationnelle ?}
{Parce que c'est une association datée, avec un historique.}
{Un attribut unique dans \code{conseiller\_type} écraserait l'historique. La
table permet plusieurs lignes et une ligne ouverte
(\code{date\_fin IS NULL}).}

\qa{6}{Pourquoi \code{ref\_parrain} est-il ajouté par \code{ALTER TYPE} ?}
{Pour démontrer l'évolution d'un type et respecter le parrainage demandé.}
{Recréer \code{client\_type} aurait cassé les dépendances.
\code{ALTER TYPE ... ADD ATTRIBUTE ... CASCADE} ajoute l'attribut sans tout
recréer.}

\qa{7}{Pourquoi la table \code{client} a-t-elle un OID système et non
\code{OBJECT IDENTIFIER IS PRIMARY KEY} ?}
{Parce qu'un \code{REF} basé sur la PK ne peut pas être inséré dans une
\emph{nested table} (\code{ORA-22979}).}
{Les signataires stockent \code{REF client\_type} dans une \emph{nested table}.
Avec un OID système, l'insertion fonctionne.}

\qa{8}{Pourquoi la table \code{agence} a-t-elle
\code{OBJECT IDENTIFIER IS PRIMARY KEY} ?}
{Pour pouvoir utiliser \code{MAKE\_REF(agence, a.code)} dans la vue objet.}
{\code{MAKE\_REF} a besoin d'un OID basé sur une clé connue. C'est ce qui
permet de construire les collections de la vue.}

\subsection{Partie A}

\qa{9}{Combien de types, tables et contraintes ?}
{19 types métier, 11 tables logiques plus une table de stockage de
\emph{nested table}, 4 séquences, 13 FK, 6 triggers et 9 sous-programmes
catalogués.}
{\code{USER\_OBJECTS} affiche 26 lignes de type \code{TYPE} : il y a des types
internes \code{SYS\_YOID...} et une version supplémentaire de
\code{CLIENT\_TYPE} due à l'évolution. Les 19 types métier sont dans
\code{USER\_TYPES}.}

\qa{10}{Comment sont stockés les téléphones ?}
{Dans un \code{VARRAY(3) OF VARCHAR2(20)}, inline dans la table
\code{client}.}
{On peut les interroger avec \code{TABLE(c.telephones)}.
\code{CARDINALITY} ne fonctionne pas sur un \code{VARRAY} dans cette version
(\code{ORA-00932}).}

\qa{11}{Comment sont stockés les signataires ?}
{Dans une \emph{nested table} \code{signataire\allowbreak\_nt\_type}, stockée dans
\code{signataire\allowbreak\_nt\_st} via \code{NESTED TABLE ... STORE AS}.}
{Chaque élément est un \code{signataire\_elt\allowbreak\_nt\_type} avec
\code{ref\_client} et \code{droit}. La contrainte
\code{droit IN ('X','R','D')} est ajoutée sur la table de stockage par
\code{ALTER TABLE}.}

\qa{12}{Pourquoi \code{AUDIT} est-il entre guillemets ?}
{Parce que \code{AUDIT} est un mot réservé Oracle.}
{Sans guillemets, \code{CREATE TABLE AUDIT} échoue. On écrit \code{"AUDIT"}
partout.}

\qa{13}{Pourquoi la contrainte \code{ref\_client IS NOT NULL} de
\code{cpt\_courant} est-elle différée ?}
{Pour permettre B5b : insérer un compte courant sans titulaire, puis le mettre
à jour avant le \code{COMMIT}.}
{Une contrainte immédiate refuserait l'insertion. La contrainte
\code{ck\_cc\_client} est \code{DEFERRABLE INITIALLY DEFERRED} et vérifiée au
moment de la validation. Les contraintes équivalentes de \code{cpt\_epargne}
et \code{cpt\_terme} restent immédiates.}

\qa{14}{Comment fonctionne l'intégrité des \code{REF} ?}
{Oracle vérifie les FK sur les colonnes et attributs \code{REF}.}
{Un \code{REF} est un pointeur ; la FK garantit que l'objet référencé existe.
Pour supprimer un client, il faut supprimer ou mettre à \code{NULL} les
références.}

\qa{15}{Pourquoi pas de FK ni de trigger sur la \emph{nested table} des
signataires ?}
{Oracle l'interdit sur le stockage d'une \emph{nested table}
(\code{ORA-30730}, \code{ORA-22812}).}
{On contrôle les droits dans le trigger sur \code{operation} et dans les
méthodes (\code{estSignataire}, \code{peutRetirer}).}

\qa{16}{Comment le script est-il rejouable ?}
{Il supprime la vue, les tables, les séquences et les types dans le bon ordre,
puis recrée tout.}
{\code{purge\_donnees.sql} fait des \code{DELETE} pour vider les données sans
supprimer le schéma. Le script unique est une concaténation A à E.}

\subsection{Partie B}

\qa{17}{Que démontre B5b ?}
{Une contrainte \code{REF IS NOT NULL} différée sur \code{cpt\_courant} : on
insère CC99 sans titulaire, puis on met à jour avant le \code{COMMIT}.}
{C'est un cas typique où la vérification doit attendre la fin de la
transaction. Les comptes d'épargne et à terme gardent une contrainte
immédiate.}

\qa{18}{Comment \code{change\_Portable} gère-t-il les \code{VARRAY} ?}
{Elle vérifie \code{COUNT} : si un élément, \code{EXTEND} ; si deux ou trois,
remplacement ; si \code{NULL} ou vide, recréation avec deux éléments.}
{Elle gère aussi \code{NO\_DATA\_FOUND} pour un client inexistant. Elle fait
\code{COMMIT} en succès et \code{ROLLBACK} en erreur.}

\qa{19}{Que fait \code{drop\_client} ?}
{Elle supprime les opérations, retire le client des listes de signataires, met
à \code{NULL} les parrainages, supprime ses comptes, puis supprime le client.}
{Elle fait \code{COMMIT} en succès et \code{ROLLBACK} en erreur. C'est une
opération de gestion autonome.}

\qa{20}{Pourquoi \code{drop\_client} fait \code{COMMIT} et \code{virement}
non ?}
{Parce que \code{drop\_client} est une opération finale, alors que
\code{virement} est une brique réutilisable.}
{Le \code{COMMIT} de \code{virement} empêcherait l'appelant de rollback
l'ensemble. \code{drop\_client} doit persister.}

\subsection{Partie C}

\qa{21}{Comment accéder à \code{c.adresse.ville} ?}
{Avec un alias de table : \code{SELECT c.adresse.ville FROM client c}.}
{Sans alias, Oracle ne sait pas déréférencer et lève \code{ORA-00904}. Pour le
parrain du parrain : \code{c.ref\_parrain.ref\_parrain.nom}.}

\qa{22}{Comment interroger un \code{VARRAY} et une \emph{nested table} ?}
{Avec \code{TABLE(...)}.}
{\code{TABLE(c.telephones)} pour le \code{VARRAY},
\code{TABLE(cc.signataires)} pour la \emph{nested table}. On peut joindre avec
\code{s.ref\_client = REF(c)}.}

\qa{23}{Pourquoi \code{CARDINALITY} échoue-t-il sur un \code{VARRAY} ?}
{Parce que cette version d'Oracle ne l'accepte pas sur un \code{VARRAY}
(\code{ORA-00932}).}
{On utilise \code{TABLE(...)} et \code{COUNT(*)}. \code{CARDINALITY}
fonctionne sur les \emph{nested tables}.}

\qa{24}{Comment fonctionne 6.20 ?}
{On ouvre la \emph{nested table} des signataires de CC7 et on joint au client
par \code{s.ref\_client = REF(c)}.}
{On affiche l'adresse complète en concaténant \code{c.adresse.num\_rue},
\code{c.adresse.rue} et \code{c.adresse.ville}. Résultat : clients 2 (droit D),
3 (droit X) et 5 (droit R).}

\qa{25}{Que fait 6.25 ?}
{Elle supprime les comptes à terme expirés depuis plus d'un an et dont le
montant bloqué est nul.}
{CT4 est le seul concerné ; CT1, CT2 et CT3 ont un montant bloqué supérieur à
0 et sont conservés. La requête fait un \code{COMMIT}.}

\subsection{Partie D}

\qa{26}{Différence entre méthode membre, méthode statique et procédure
cataloguée ?}
{Membre : sur une instance ; statique : sur le type ; cataloguée : autonome.}
{\code{fermeCompte} est statique car elle prend un numéro de compte, pas une
instance. \code{virement} est cataloguée car elle manipule plusieurs comptes.}

\qa{27}{Pourquoi \code{fraisTenue} est-elle déclarée avec \code{OVERRIDING}
dès le \code{CREATE TYPE} ?}
{Parce qu'\code{ALTER TYPE ADD MEMBER} ne permet pas de créer un override
(\code{ORA-22324} / \code{PLS-00636}).}
{On déclare la méthode dans la super-classe, puis on écrit
\code{OVERRIDING MEMBER FUNCTION fraisTenue} dans les corps des sous-types.}

\qa{28}{Pourquoi \code{passeMvt} ne fait-elle pas \code{COMMIT} ?}
{Pour que l'appelant contrôle la transaction.}
{Elle insère l'opération, le trigger \code{AFTER} met à jour le solde, puis
elle recalcule le solde depuis \code{operation} : c'est idempotent. Un
\code{COMMIT} interne empêcherait les tests de rollback.}

\qa{29}{Comment \code{virement} garantit-il l'atomicité ?}
{Avec \code{SAVEPOINT} et \code{ROLLBACK TO} en cas d'erreur.}
{Les deux insertions sont dans la même transaction. La commission
inter-agences est de 1~\% avec un minimum de 100~DA. Le test CC3 vers CC4
échoue sur le découvert et CC4 reste inchangé.}

\qa{30}{Comment \code{fermeCompte} gère-t-elle les erreurs ?}
{\code{NO\_DATA\_FOUND} pour un compte inexistant ; \code{WHEN OTHERS} avec
\code{ROLLBACK}.}
{Elle cherche dans CC, CE et CT, supprime les opérations si c'est un CC,
supprime le compte, puis \code{COMMIT}. Elle est testée en dernier car elle
est destructive.}

\qa{31}{Pourquoi \code{trg\_op\_avant} est-il un \code{COMPOUND TRIGGER} ?}
{Pour éviter \code{ORA-04091} en lisant \code{operation} pour le plafond
journalier.}
{\code{BEFORE STATEMENT} prend l'instantané des retraits du jour ;
\code{BEFORE EACH ROW} accumule. On ne relit jamais la table mutante dans la
partie ligne par ligne.}

\qa{32}{Quel est l'ordre des triggers et comment le solde est-il mis à
jour ?}
{\code{BEFORE} valide, \code{AFTER INSERT} met à jour le solde.}
{\code{trg\_op\_avant} refuse les opérations invalides ;
\code{trg\_op\_apres} applique \code{solde = solde + montant - commission}.
\code{passeMvt} recalcule ensuite depuis \code{operation}, sans double
comptage.}

\qa{33}{Comment sont calculés les intérêts et la sortie anticipée ?}
{Épargne : \code{solde * (1 + taux/100)}. Terme :
\code{solde * (1 + taux/100 * durée/12)}. Sortie anticipée : 95~\% du solde.}
{CE5 : $500\,000 \times 1,03 = 515\,000$ ; sortie 475\,000. CT1 :
$200\,000 \times 1,032 = 206\,400$ ; sortie 190\,000.
\code{afficheInterets} affiche les intérêts seuls : 3\,750 + 15\,000 =
18\,750 pour le client 1.}

\qa{34}{Comment fonctionne \code{maj\_statuts\_dormants} ?}
{Elle passe en \code{DORMANT} les comptes courants \code{ACTIF} sans opération
depuis douze mois.}
{Le test est fait dans un bloc avec \code{ROLLBACK} pour ne pas modifier
l'état final. Le résultat affiché est « 1 compte(s) passe(s) DORMANT ».}

\qa{35}{Comment \code{scan\_echeances} évite-t-elle les doublons ?}
{Avec un test \code{COUNT(*) = 0} sur \code{alerte} pour le même compte et le
même type, avec \code{traitee = 'N'}.}
{Le trigger \code{trg\_ct\_echeance} crée les alertes à l'insertion ou à la
mise à jour ; \code{scan\_echeances} fait un balayage global pour les comptes
existants.}

\subsection{Partie E}

\qa{36}{Expliquez le type récursif et la double notation.}
{\code{ref\_parrain REF client\_type} permet de chaîner les parrains.
\code{c.ref\_parrain.ref\_parrain.nom} donne le grand-parrain.}
{E10 met à jour les parrains des clients 4, 5 et 6. Les clients 2 et 3 restent
parrainés par 1. La requête affiche parrain et grand-parrain.}

\qa{37}{Expliquez la collection multi-niveaux \code{releves}.}
{\code{releves} contient un \code{VARRAY(12)} de mois ; chaque mois contient
une \emph{nested table} de statistiques par type.}
{On interroge avec un double \code{TABLE} :
\code{TABLE(r.mois\_vry) m, TABLE(m.types) t}. Le relevé 2026 a cinq mois. Les
totaux sont DEPOT 13/880\,000, RETRAIT 32/620\,000, VIREMENT 3/130\,000,
CB 2/35\,000.}

\qa{38}{Expliquez la vue objet \code{Agence\_VOR}.}
{C'est une vue de type \code{agence\_vor\_type} avec deux collections calculées
par \code{CAST(MULTISET(...))}.}
{\code{WITH OBJECT IDENTIFIER (code)} donne un OID à la vue.

\code{MAKE\_REF(agence,\allowbreak\ a.code)} fabrique les \code{REF} d'agence.
On interroge les collections avec \code{TABLE(v.\allowbreak conseillers)} et
\code{TABLE(v.\allowbreak comptes)}.}

\qa{39}{Expliquez la méthode \code{mutation} et le problème de
\code{date\_fin}.}
{Elle ferme l'affectation courante et en ouvre une nouvelle. \code{date\_fin}
doit être strictement supérieure à \code{date\_debut}.}
{Si on mute le même jour, \code{SYSDATE} peut être égal à \code{date\_debut},
ce qui viole \code{CK\_AFF\_DATES} (\code{ORA-02290}). On utilise
\code{GREATEST(SYSDATE, date\_debut) + 1/86400}.}

\qa{40}{Pourquoi la mutation ne remplace-t-elle pas automatiquement le
responsable ?}
{Parce que le sujet demande de signaler, pas de remplacer.}
{Le responsable est un rôle de management. Le remplacer automatiquement
pourrait violer les règles internes. La méthode ferme l'affectation et laisse
la décision au gestionnaire.}

\subsection{Validation et soutenance}

\qa{41}{Pourquoi y a-t-il des erreurs ORA dans la sortie ?}
{Ce sont des tests d'exception volontaires.}
{\code{ORA-20001} à \code{ORA-20004} pour les triggers,
\code{ORA-20010}/\code{ORA-20011} pour la mutation, et les
\code{ORA-06512}/\code{ORA-04088} associés. Les cinq tests de triggers sont
capturés par \code{EXCEPTION WHEN OTHERS} et affichés comme « Refus attendu » ;
les tests \code{passeMvt}, \code{virement} et \code{mutation} laissent
volontairement l'erreur s'afficher pour prouver le refus. Aucune erreur de
compilation.}

\qa{42}{Pourquoi zéro objet invalide est-il important ?}
{Parce qu'après \code{ALTER TYPE ... CASCADE}, les objets dépendants peuvent
devenir \code{INVALID}.}
{On recompile toutes les procédures et fonctions à la fin de la partie D. Le
contrôle final \code{WHERE status <> 'VALID'} ne retourne aucune ligne.}

\qa{43}{Que se passe-t-il si on relance le script ?}
{Il supprime et recrée tout ; le résultat final est identique.}
{Les séquences sont recréées, les données réinsérées. Ce n'est pas un script
de production, c'est un script de régénération de TP.}

\qa{44}{Pourquoi le script unique est-il une concaténation et pas un
\code{@} ?}
{Pour qu'il soit portable et exécutable seul, sans dépendre du chemin des
parties.}
{Les parties restent séparées pour la lisibilité et l'évaluation. Le script
unique ajoute des séparateurs et un contrôle final.}

\qa{45}{Quelles sont les limites de votre projet ?}
{Opérations uniquement sur les comptes courants ; virement CC vers CC ; pas de
remplacement automatique du responsable ; pas d'interface graphique ; pas de
tests de concurrence.}
{Ces limites sont documentées dans \code{ETAPE\_0\_CADRAGE.md}. On peut les
justifier par le sujet ou par le temps. Par exemple,
\code{operation\_type.ref\_compte} est \code{REF cptCourant\_type} : un
virement vers l'épargne demanderait un modèle polymorphe.}

\section{Les quinze pièges à éviter}

\begin{enumerate}
  \item Ne pas confondre \code{CARDINALITY} et \code{COUNT} :
        \code{CARDINALITY} sur \emph{nested table}, \code{TABLE} +
        \code{COUNT} sur \code{VARRAY}.
  \item Ne pas oublier l'alias pour \code{c.adresse.ville},
        \code{c.ref\_conseiller.nom}, \code{c.ref\_parrain.nom}.
  \item Ne pas dire que les ORA sont des bugs : ce sont des tests.
  \item Ne pas dire que \code{virement} fait \code{COMMIT} : il ne le fait
        pas.
  \item Ne pas dire que \code{passeMvt} double le solde : le recalcul est
        idempotent.
  \item Ne pas confondre \code{trg\_op\_avant} et \code{trg\_op\_apres} :
        avant = validation, après = solde et alertes.
  \item Ne pas oublier \code{GREATEST(SYSDATE, date\_debut) + 1/86400} dans
        la mutation.
  \item Ne pas dire que \code{client} a un OID basé sur la PK : c'est
        \code{agence}.
  \item Ne pas dire que \code{AUDIT} est une table normale : il est entre
        guillemets.
  \item Ne pas dire que la \emph{nested table} des signataires a des FK ou des
        triggers : Oracle l'interdit.
  \item Ne pas capturer la partie C après la partie D : CC1, CE4 et
        l'opération de 1,5~M ne sont plus là.
  \item Ne pas oublier que CT4 est supprimé par 6.25.
  \item Ne pas oublier que CE4 et CC1 sont supprimés par \code{fermeCompte}.
  \item Ne pas dire que la vue objet est une table : c'est une vue de type
        objet.
  \item Ne pas promettre une fonctionnalité non implémentée : dire « limite
        documentée » et expliquer comment on la ferait.
\end{enumerate}

\section{Phrases de secours}

Quelques formulations permettent de garder le contrôle lorsqu'une question
sort du périmètre préparé.

\begin{itemize}
  \item « Je n'ai pas implémenté ce point, mais voici comment je le ferais… »
  \item « Ce n'est pas un bug : c'est un test d'exception volontaire, voici le
        bloc qui le capture. »
  \item « La sortie complète est dans \code{\_FULL.out}, je peux vous montrer
        la section exacte. »
  \item « Ce choix est justifié dans \code{ETAPE\_0\_CADRAGE.md}, hypothèse
        H… »
  \item « Je préfère ne pas modifier l'état en direct ; je vous montre la
        sortie déjà exécutée. »
\end{itemize}

\section{Chiffres à retenir}

\begin{table}[htbp]
  \centering
  \small
  \caption{Chiffres clés du projet à la fin du rejeu complet.}
  \label{tab:chiffres-cles}
  \begin{tabularx}{\textwidth}{@{}L{7.0cm}C{3.0cm}X@{}}
    \toprule
    \entete{Élément} & \entete{Valeur} & \entete{Commentaire} \\
    \midrule
    Clients & 6 & \\
    Conseillers & 4 & \\
    Agences & 3 & \\
    Affectations finales & 9 & 5 initiales + 4 mutations. \\
    Comptes courants finaux & 6 & CC1 supprimé. \\
    Comptes d'épargne finaux & 4 & CE4 supprimé. \\
    Comptes à terme finaux & 3 & CT4 supprimé par 6.25. \\
    Opérations finales & 15 & 20 au départ, 5 supprimées avec CC1. \\
    Signataires finaux & 9 & \\
    Alertes & 2 & CT2 et CT3. \\
    Audit & 0 & Lignes annulées par rollback. \\
    Types métier & 19 & \\
    Tables logiques & 11 + 1 stockage & \\
    Séquences & 4 & \\
    Triggers & 6 & \\
    Sous-programmes catalogués & 9 & \\
    Objets invalides & 0 & \\
    \bottomrule
  \end{tabularx}
\end{table}

Les soldes finaux des comptes courants sont : CC2 $-50\,000$ ; CC3
38\,000~DA ; CC4 325\,000~DA ; CC5 0~DA et \code{BLOQUE} ; CC6 40\,000~DA ;
CC7 185\,000~DA. Les comptes d'épargne finaux sont CE1 150\,000~DA, CE2
250\,000~DA, CE3 80\,000~DA et CE5 500\,000~DA. Les comptes à terme finaux
sont CT1 200\,000~DA, CT2 100\,000~DA et CT3 50\,000~DA.

\begin{attentionbox}[Un chiffre, un moment]
Les 20 opérations et les 7 comptes courants sont des valeurs du début de la
partie C. Les 15 opérations et les 6 comptes courants sont des valeurs de fin
de partie D. Les deux sont exacts ; il faut simplement préciser le moment
auquel on se réfère.
\end{attentionbox}

\section{Checklist avant la soutenance}

\begin{itemize}
  \item La base répond et \code{MASTER1} existe.
  \item \code{projet\_bdda\_banque.sql} a été exécuté en entier (code de sortie
        0).
  \item \code{SELECT COUNT(*) FROM user\_objects WHERE status <> 'VALID';}
        retourne 0.
  \item Les cinq parties et les fichiers correspondants sont connus.
  \item Trois choix sont expliqués : héritage, OID client, trigger
        \code{COMPOUND}.
  \item Trois preuves sont montrées : \code{nbCepargne}, \code{passeMvt},
        \code{Agence\_VOR}.
  \item Trois erreurs attendues sont expliquées : \code{ORA-20001},
        \code{ORA-20003}, \code{ORA-20010}.
  \item La disparition de CC1, CE4 et CT4 est expliquée.
  \item Les captures d'écran sont prêtes, conformément à
        l'annexe~\ref{ann:captures}.
  \item \code{\_FULL.out} est disponible en secours.
\end{itemize}

\section{Conclusion de la validation}

La validation ne consiste pas à réciter quarante-cinq réponses. Elle consiste
à montrer un projet cohérent, à expliquer les choix et à distinguer clairement
les erreurs volontaires des erreurs de compilation. Les preuves existent : le
script complet a été exécuté, les objets sont valides et les résultats sont
reproductibles. La soutenance peut donc se concentrer sur la démonstration et
sur la justification, plutôt que sur la réparation.
